\section{Introduction}%
\label{sec:introduction}

We have a sequence of data \((X_i)_{i=1}^{137}\) from a clinical experiment of a veteran disease trial with the columns \verb|treat|, \verb|cell|,
\verb|survt|, \verb|cens|, \verb|KAR|, \verb|didur|, \verb|age| and \verb|prith|.

\begin{itemize}
    \item[\texttt{treat}] whether the patient was treated (\(1\)) or received placebo (\(2\))
    \item[\texttt{cell}] cell type of the patient (squamous (\(1\)), small (\(2\)), adeno (\(3\)), large (\(4\)))
    \item[\texttt{survt}] survival time (days)
    \item[\texttt{cens}] status (censored (\(1\)), died (\(0\)))
    \item[\texttt{KAR}]  performance status - Karnofsky score (worst (\(0\)),\dots,best (\(100\)))
    \item[\texttt{didur}] disease duration from diagnosis to treatment (months)
    \item[\texttt{age}]  age of the patient (years)
    \item[\texttt{prith}] prior therapy received by the patient (none (\(0\)), some (\(10\)))
\end{itemize}

I was randomly assigned the number \emph{\Romannumeral{26}}, which corresponds to the selection

\begin{codesample}
disease duration didur<=9 (months) versus didur>9 (months) při treat=2
\end{codesample}

In agreement with this, I will denote \((A_i)_{i=1}^{48}\) the conditioned data \((X \in (X_i) \mid X_{\verb|didur|} \le 9, \,
X_{\verb|treat|} = 2)_i\) and \((B_i)_{i=1}^{20}\) the conditioned data \((X \in (X_i) \mid X_{\verb|didur|} > 9, \,
X_{\verb|treat|} = 2)_i\). We will also denote \((A^{(1)}_i) = (A_i \mid \delta(i) = 1)\) and similarly 
\((B^{(1)}_i) = (B_i \mid \delta(i) = 1)\) and analogously \((A^{(0)}_i)\) and 
\((B^{(0)}_i)\) for the censored data.


First of all, let us see some elementary statistics of \(A\) and \(B\).

\begin{table}[h]
    \centering
    \begin{tabularx}{\linewidth}{r>{\centering\arraybackslash}X>{\centering\arraybackslash}X>{\centering\arraybackslash}X>{\centering\arraybackslash}X>{\centering\arraybackslash}X>{\centering\arraybackslash}X>{\centering\arraybackslash}X>{\centering\arraybackslash}X>{\centering\arraybackslash}X}
\toprule  Group & $\#$ & $\sum \delta$ & $\mathop{E}\delta$ & $\overline{X^{(1)}_n}$ & $\std(X^{(1)})$ & $\overline{X^{(0)}_n}$ & $\std(X^{(0)})$ & $\min(X^{(1)})$ & $\max(X^{(1)})$\\\midrule
  $A$ & 48 & 37 & $0.77$ & $124.35$ & $197.58$ & $114.73$ & $106.93$ & 1 & 991\\\addlinespace
  $B$ & 20 & 14 & $0.7$ & $172.5$ & $266.67$ & $73.33$ & $70.47$ & 1 & 999\\\bottomrule
\end{tabularx}
    \caption{Elementary statistics of the data groups. \(\mathop{E} \delta \) denotes the proportion of
    non-censored observations in the set.}
    \label{tab:elementary_stats}
\end{table}

% TODO: Add censored data histogram formula and reimplement the histogram to use it <23-09-23> 
\begin{figure}
    \centering
    \incfig{histograms}
    \caption{Histograms of the censored lifetimes.}
    \label{fig:histograms}
\end{figure}

\section{Parametric Estimates}%
\label{sec:parametric_estimates}

Based on the histograms in \Cref{fig:histograms}, we will first attempt to fit an exponential distribution to the lifetime data.

\begin{theorem}[Likelihood functions for the \ac{RC} model]
    \label{thm:likelihood}
    Let \(F_T=F_T(t,\bs{\theta}_1)\) and \(F_C = F_C(c, \bs{\theta}_2)\) be distributions of \(T\) and \(C\)
    respectively, such that \(F_T\) and \(F_C\) do not contain the same parameters (even in a relationship sense). Let
    us have \((W_j, \delta_j)_1^n\) \ac{iid} \(F_{W,\delta}\) replications. Then
    \[
        L(\bs{\theta}_1, \bs{\theta}_2) = f_{\bs{W}, \bs{\delta}}(\bs{w}, \bs{\delta}; \bs{\theta}_1, \bs{\theta}_2) =
        L_1(\bs{\theta}_1) \cdot L_2(\bs{\theta}_2),
    \] where \[
        \begin{aligned}
            L_1(\bs{\theta}_1) &= \prod_{j\in D} f_T(w_j) \prod_{j \in V} (1 - F_T(w_j)) \\
            L_2(\bs{\theta}_2) &= \prod_{j\in D} (1-F_C(w_j)) \prod_{j \in V} f_C(w_j),
        \end{aligned}
    \] with the notation \(D = \left\{ j \mid \delta_j = 1 \right\}\) and \(V = \left\{ j\mid \delta_j = 0 \right\}\).
\end{theorem}

Using \Cref{thm:likelihood}, we can make a \ac{MLE} of the distribution \(F_T\) without assuming a distribution
of \(C\). Specifically for a \ac{MLE} of an exponentially distributed \(F_T\), we have \[
    L_1(\lambda) = \prod_{j\in D} f_T(w_j) \cdot \prod_{j \in V} R_T(w_j) = \prod_{j\in D} \lambda e^{-\lambda w_j}
    \cdot \prod_{j \in V} e^{-\lambda w_j} = \lambda^{|D|} \cdot \exp(-\lambda \sum_{j = 1}^n w_j),
\] and the \ac{LEs} and \ac{MLE} estimate follows to be
\begin{gather*}
    \ell_1(\lambda) = \ln L_1(\lambda) = |D| \cdot \ln \lambda - \lambda \sum_{j = 1}^n w_j \\
    \partial_\lambda\ell_1(\lambda) = \frac{|D|}{\lambda} - \sum_{j = 1}^n w_j = 0 \implies \hat{\lambda}_{\text{ML}}
        = \frac{|D|}{\sum_{j = 1}^n w_j}.
\end{gather*}

For our data the \ac{MLE} equate to the results in \Cref{tab:MLE_exponential}

\begin{table}[h]
    \centering
    \begin{tabularx}{\linewidth}{r>{\centering\arraybackslash}X}
\toprule  $Group$ & $\hat{\lambda}_\text{ML}$\\\midrule
  $A$ & $0.0082$\\\addlinespace
  $B$ & $0.007$\\\bottomrule
\end{tabularx}
    \caption{\ac{MLE} estimates for the exponential distribution of \(F_T\).}
    \label{tab:MLE_exponential}
\end{table}

Checking the fit via a histogram while keeping the information about the data censoring is done via the \ac{KM}
estimate of the \ac{CDF}. The \ac{KM} model for censored data gives us the estimate for the survival function as
\[
    \widehat{S}(t)=\prod_{i:t_{i}\leq t}\left(1-{\frac{d_{i}}{n_{i}}}\right).
\] 

Therefore an estimate of the \ac{CDF} that takes the censors into account is 
\[
    \widehat{F}(t)=1-\widehat{S}(t).
\]

We can then create a Q-Q plot that accounts for the censors simply by evaluating the quantile at the
estimated quantile from the \ac{KM} model for each time and plotting against the true observed time. This equates to the
plot being \[
    \left\{ F^{-1}_T(\widehat{F}(t)), t  \right\} \qcomma  \forall t \in (t_i),
\]  where \(F_T\) is the distribution of that we want to check.

A P-P plot accounting for the censors is analogously \[
    \left\{ F_T(t), \widehat{F}(t) \right\} \qcomma  \forall t \in (t_i).
\] 


\begin{figure}[h]
    \centering
    \begin{subfigure}{0.4\textwidth}
        \incfig{histogram_vs_parametric_A}
        \subcaption{Comparison of histogram and \ac{MLE} estimates.}
    \end{subfigure} 
    \begin{subfigure}{0.4\textwidth}
        \incfig{QQ_exponential_A}
        \subcaption{Q-Q plot of the fitted exponential \ac{MLE}}
    \end{subfigure} 
    \begin{subfigure}{0.4\textwidth}
        \incfig{PP_exponential_A}
        \subcaption{P-P plot of the fitted exponential \ac{MLE} against the \ac{KM} probability estimate.}
    \end{subfigure}
    \caption{\ac{MLE} fit of an exponential distribution of group \(A\) justification plots.}
    \label{fig:histogram_vs_parametric_and_QQ_A}
\end{figure}

\begin{figure}[h]
    \centering
    \begin{subfigure}{0.4\textwidth}
        \incfig{histogram_vs_parametric_B}
        \subcaption{Comparison of histogram and \ac{MLE} estimates.}
    \end{subfigure} 
    \begin{subfigure}{0.4\textwidth}
        \incfig{QQ_exponential_B}
        \subcaption{Q-Q plot of the fitted exponential \ac{MLE}}
    \end{subfigure} 
    \begin{subfigure}{0.4\textwidth}
        \incfig{PP_exponential_B}
        \subcaption{P-P plot of the fitted exponential \ac{MLE} against the \ac{KM} probability estimate.}
    \end{subfigure}
    \caption{\ac{MLE} fit of an exponential distribution of group \(B\) justification plots.}
    \label{fig:histogram_vs_parametric_and_QQ_B}
\end{figure}

From \Cref{fig:histogram_vs_parametric_and_QQ_A} and \Cref{fig:histogram_vs_parametric_and_QQ_B}, we see that
the exponential distribution is a not an ideal fit for neither of groups \(A\) and \(B\) since they seem to have heavier
tails than the fitted distribution, but is acceptable for a simple model.

Lets also try the Weibull distribution.

Similarly to the exponential distribution, we determine for the Weibull distribution
\begin{gather*}
    R_T(t) = \exp{-(\lambda t)^\alpha} = \left\{ \lambda^\alpha \eqqcolon \lambda' \right\} = \exp{-\lambda't^\alpha} \\
    f_T(t) = \lambda' \alpha t^{(\alpha - 1)} \exp{-\lambda't^\alpha}; \quad \lambda' > 0,\, \alpha > 0,\, t > 0
\end{gather*} the log-likelihood as
\[
    \ell_1(\alpha, \lambda') = |D| \ln \lambda' + |D| \ln \alpha + (\alpha - 1) \ln\prod_{j \in D} w_j - \lambda'
    \sum_{j = 1}^n w_j^\alpha.
\] % and \ac{LEs} and the \ac{MLE} to be
% \begin{gather}
%     \partial_{\lambda'} \ell_1 = |D| / \lambda' - \sum_{j = 1}^n w_j^\alpha = 0 \\
%     \partial_\alpha \ell_1 = |D|/\alpha + \sum_{n = 0}^{N} 
% \end{gather}

\newpage
\pagebreak
\warningbox{The fits of the Weibull distributions seem to be wrong. This is the code that is used to generate them.
    \begin{center}
        \includegraphics[width=0.8\textwidth]{Weibull_loglike.png}
    \end{center}
}

\begin{figure}[h]
    \centering
    \begin{subfigure}{0.4\textwidth}
        \incfig{histogram_vs_parametric_weib_A}
        \subcaption{Comparison of histogram and \ac{MLE} estimates.}
    \end{subfigure} 
    \begin{subfigure}{0.4\textwidth}
        \incfig{QQ_weib_A}
        \subcaption{Q-Q plot of the fitted Weibull \ac{MLE}}
    \end{subfigure} 
    \begin{subfigure}{0.4\textwidth}
        \incfig{PP_weibull_A}
        \subcaption{P-P plot of the fitted Weibull \ac{MLE} against the \ac{KM} probability estimate.}
    \end{subfigure}
    \caption{\ac{MLE} fit of an Weibull distribution and a QQ plot of group \(A\).}
    \label{fig:histogram_vs_parametric_and_QQ_weib_A}
\end{figure}

\begin{figure}[h]
    \centering
    \begin{subfigure}{0.4\textwidth}
        \incfig{histogram_vs_parametric_weib_B}
        \subcaption{Comparison of histogram and \ac{MLE} estimates.}
    \end{subfigure} 
    \begin{subfigure}{0.4\textwidth}
        \incfig{QQ_weib_B}
        \subcaption{Q-Q plot of the fitted Weibull \ac{MLE}}
    \end{subfigure} 
    \begin{subfigure}{0.4\textwidth}
        \incfig{PP_weibull_B}
        \subcaption{P-P plot of the fitted Weibull \ac{MLE} against the \ac{KM} probability estimate.}
    \end{subfigure}
    \caption{\ac{MLE} fit of an Weibull distribution of group \(B\) justification plots.}
    \label{fig:histogram_vs_parametric_and_QQ_weib_B}
\end{figure}

\begin{figure}[h]
    \centering
    \begin{subfigure}[b]{0.45\textwidth}
        \incfig[1.1\textwidth]{log_likelihood_surf_A}
        \subcaption{Log-likelihood surface in the neighbourhood of the \ac{MLE} estimate of A}
        \label{fig:log_likelihood_surf_A}
    \end{subfigure}
    \hfill
    \begin{subfigure}[b]{0.45\textwidth}
        \incfig[1.1\textwidth]{log_likelihood_surf_B}
        \subcaption{Log-likelihood surface in the neighbourhood of the \ac{MLE} estimate of B}
        \label{fig:log_likelihood_surf_B}
    \end{subfigure}
\end{figure}

As my parametric fit, I will use the \ac{MLE} of the exponential for lack of a better fit with Weibull.

\section{\ac{KM} Plots}%
\label{sec:KM_plot}

\begin{figure}[h]
    \centering
    \begin{subfigure}[b]{0.45\textwidth}
        \incfig{KM_A}
        \subcaption{\ac{KM} plot of group A against the \(R_\text{MLE}\) of the exponential distribution.}
        \label{fig:KM_A}
    \end{subfigure}
    \hfill
    \begin{subfigure}[b]{0.45\textwidth}
        \incfig{KM_B}
        \subcaption{\ac{KM} plot of group B against the \(R_\text{MLE}\) of the exponential distribution.}
        \label{fig:KM_B}
    \end{subfigure}
\end{figure}

\begin{figure}[h]
    \incfig{KM_AB}
    \subcaption{\ac{KM} plot of group A and group B.}
    \label{fig:KM_AB}
\end{figure}


\section{\ac{NA} Plots}%
\label{sec:NA_plot}

\begin{figure}[h]
    \centering
    \begin{subfigure}[b]{0.45\textwidth}
        \incfig{NA_A}
        \subcaption{\ac{NA} plot of group A against the \(\Lambda_\text{MLE}\) of the exponential distribution.}
        \label{fig:NA_A}
    \end{subfigure}
    \hfill
    \begin{subfigure}[b]{0.45\textwidth}
        \incfig{NA_B}
        \subcaption{\ac{NA} plot of group B against the \(\Lambda_\text{MLE}\) of the exponential distribution.}
        \label{fig:NA_B}
    \end{subfigure}
\end{figure}

Based on the \ac{NA} plots in \cref{fig:NA_A} and \cref{fig:NA_B}, both \(A\) and \(B\) have a \ac{DFR}.

\section{Comparison of \(A\) and \(B\)}%
\label{sec:comparison}

As we are comparing the exponentially distributed model of lifetimes, we already know how the comparison will look for
\(R_A\) and \(R_B\) linear \(\Lambda_A\) and \(\Lambda_B\) and the constant \(\lambda_A\) and \(\lambda_B\), but for completeness, here are the comparison graphs.

\begin{figure}[h]
    \centering
    \begin{subfigure}[b]{0.45\textwidth}
        \incfig{comp_R}
        \subcaption{Comparison of \(R_A\) and \(R_B\).}
        \label{fig:comp_R}
    \end{subfigure}
    \hfill
    \begin{subfigure}[b]{0.45\textwidth}
        \incfig{comp_λ}
        \subcaption{Comparison of \(\lambda_A\) and \(\lambda_B\).}
        \label{fig:comp_λ}
    \end{subfigure}

    \begin{subfigure}[b]{0.45\textwidth}
        \incfig{comp_Λ}
        \subcaption{Comparison of \(\Lambda_A\) and \(\Lambda_B\).}
        \label{fig:comp_Λ}
    \end{subfigure}
    \caption{Comparisons of descriptive model functions.}
\end{figure}

The \ac{MTTF} of an exponential hazard rate is \(1/\lambda\) and the median lifetime \(t_\text{med}\) is
given by \begin{align*}
    Pr(T > t_\text{med}) &= 1/2 \\
    e^{-\lambda t_\text{med}} &= 1/2 \\
    t_\text{med} &= \ln{2}/\lambda.
\end{align*}

\begin{table}[h]
    \centering
    \begin{tabularx}{\linewidth}{r>{\centering\arraybackslash}X>{\centering\arraybackslash}X}
\toprule  Group & $MTTF$ & $t_\text{med}$\\\midrule
  A & $120.0$ & $85.0$\\\addlinespace
  B & $140.0$ & $99.0$\\\bottomrule
\end{tabularx}
    \caption{Comparison of \ac{MTTF} and the median lifetime for groups \(A\) and \(B\).}
    \label{tab:comparison}
\end{table}

\section{Log-log \(R\) Plots}%
\label{sec:loglogR_plot}

\begin{figure}[h]
    \centering
    \incfig{log_log_R}
    \caption{Log-log \(R\) plots of \(R_A\) and \(R_B\).}
    \label{fig:log_log_R}
\end{figure}

As we can see from the \Cref{fig:log_log_R}, there are several crossings of the two step functions and therefore the
proportional hazard model does not fit the groups \(A\) and \(B\).

\section{Summary}%
\label{sec:summary}
The main insight is given in the \Cref{fig:KM_AB}, where we can see that the survival function estimate is lower for
group \(A\) in the interval \((100, 400)\) days. I would say that the right tail of both groups does not include enough
observations for the estimate have a big weight in the conclusions. This indicates that if a patient had more than
9 months from the diagnosis before getting treated, his survival chances are worse. That is a good indication of the
  treatment being effective.

What is very interesting to me from the data is that the patients already survived 9 months as indicated by the disease
duration value but after getting the treatment, there are so many observations of death in the interval from 0 to 100
days. This is quite puzzling and I would deduce intuitively from this information that the treatment is causing the
death. Some more information about how the data was collected would be useful for putting this strange and pessimistic
theory into light of day.
